% Auto-gerado por C13_build_T1_calibrated.py - versao calibrada (Q>=2 fallback Q1)
\begin{table}[htbp]
\centering
\caption{Indicadores agregados de fluxo escolar, Brasil, por macroetapa e ano (PNADC longitudinal, 2012--2023). Calibra\c{c}\~ao $t$+1 em Q$\geq$2 com fallback Q1.}
\label{tab:t1_brasil}
\begin{threeparttable}
\small
\begin{tabular}{lrrrrrr}
\toprule
Ano & Promo\c{c}\~ao & Repet\^encia & Evas\~ao & Migra\c{c}\~ao EJA & Abandono$^{a}$ & N \\
    & (\%)            & (\%)        & (\%)     & (\%)               & (\%)            &   \\
\midrule
\multicolumn{7}{l}{\textit{Painel A: EF iniciais (1\textsuperscript{o}--5\textsuperscript{o} ano)}} \\
\midrule
2012 & 64.3 & 32.5 & 1.7 & 0.2 & 1.4 & 32.993 \\
2013 & 66.9 & 30.2 & 1.5 & 0.2 & 1.2 & 30.424 \\
2014 & 69.5 & 28.0 & 1.3 & 0.1 & 1.1 & 29.923 \\
2015 & 73.2 & 25.8 & 0.7 & 0.2 & 0.8 & 21.348 \\
2016 & 73.5 & 25.6 & 0.7 & 0.2 & 0.5 & 31.510 \\
2017 & 73.7 & 25.4 & 0.6 & 0.2 & 0.5 & 30.132 \\
2018 & 75.5 & 23.7 & 0.5 & 0.2 & 0.4 & 29.172 \\
2019 & 77.2 & 22.2 & 0.5 & 0.1 & 0.5 & 27.123 \\
2020$^{\dagger}$ & 64.3 & 35.1 & 0.6 & 0.0 & 0.6 & 15.986 \\
2021$^{\dagger}$ & 70.0 & 29.6 & 0.3 & 0.1 & 0.4 & 21.967 \\
2022 & 61.1 & 38.6 & 0.2 & 0.1 & 0.2 & 22.448 \\
2023 & 61.7 & 38.0 & 0.2 & 0.0 & 0.1 & 23.136 \\
\\[0.3em]
\multicolumn{7}{l}{\textit{Painel B: EF finais (6\textsuperscript{o}--9\textsuperscript{o} ano)}} \\
\midrule
2012 & 60.0 & 31.1 & 3.9 & 1.2 & 3.3 & 23.874 \\
2013 & 61.9 & 29.6 & 4.1 & 1.0 & 3.2 & 22.667 \\
2014 & 63.9 & 28.6 & 3.7 & 0.9 & 3.0 & 22.762 \\
2015 & 66.8 & 26.9 & 2.9 & 1.3 & 2.4 & 16.472 \\
2016 & 67.7 & 26.5 & 2.6 & 1.6 & 2.1 & 24.678 \\
2017 & 68.7 & 26.2 & 2.4 & 1.3 & 2.0 & 23.842 \\
2018 & 70.4 & 24.6 & 2.3 & 1.3 & 1.8 & 23.085 \\
2019 & 72.6 & 23.1 & 1.9 & 1.0 & 1.7 & 21.619 \\
2020$^{\dagger}$ & 60.1 & 37.6 & 1.2 & 0.3 & 0.6 & 13.710 \\
2021$^{\dagger}$ & 67.6 & 29.5 & 1.7 & 0.3 & 0.6 & 18.910 \\
2022 & 59.8 & 38.2 & 1.2 & 0.3 & 0.8 & 19.008 \\
2023 & 59.9 & 38.5 & 0.8 & 0.3 & 0.6 & 19.183 \\
\\[0.3em]
\multicolumn{7}{l}{\textit{Painel C: Ensino M\'edio (1\textsuperscript{o}--3\textsuperscript{o} ano)}} \\
\midrule
2012 & 42.9 & 18.6 & 26.7 & 1.4 & 11.8 & 15.827 \\
2013 & 42.8 & 18.3 & 26.9 & 1.1 & 11.6 & 15.088 \\
2014 & 45.1 & 17.6 & 26.9 & 1.1 & 11.4 & 15.498 \\
2015 & 46.0 & 17.2 & 25.9 & 1.3 & 10.6 & 11.184 \\
2016 & 45.9 & 17.1 & 26.1 & 1.7 & 9.2 & 16.738 \\
2017 & 46.2 & 16.8 & 26.3 & 2.2 & 10.6 & 15.718 \\
2018 & 47.9 & 16.1 & 25.9 & 1.8 & 9.7 & 14.939 \\
2019 & 48.8 & 15.7 & 25.8 & 1.5 & 9.8 & 14.166 \\
2020$^{\dagger}$ & 41.3 & 36.3 & 17.5 & 0.6 & 2.7 & 9.072 \\
2021$^{\dagger}$ & 44.5 & 25.3 & 23.6 & 0.9 & 6.1 & 12.388 \\
2022 & 39.7 & 35.2 & 18.8 & 0.4 & 6.6 & 12.666 \\
2023 & 41.3 & 33.8 & 18.1 & 0.4 & 5.8 & 12.484 \\
\bottomrule
\end{tabular}
\begin{tablenotes}\footnotesize
\item \textit{Fonte:} Painel longitudinal da PNADC trimestral (IBGE), 2012Q1--2024Q4.
\item \textit{Notas:} As taxas de promo\c{c}\~ao, repet\^encia, evas\~ao e migra\c{c}\~ao para EJA s\~ao \emph{entre-anos} (do ano $t$ ao ano $t+1$). A obs.\ de refer\^encia em $t$ e em $t+1$ \'e, por crit\'erio, o primeiro trimestre Q$\geq$2 observado no ano correspondente; se o domic\'ilio aparece apenas no Q1, esta obs.\ \'e usada (Se\c{c}\~ao~\ref{ssec:efeito_ferias}). Abandono \'e \emph{intra-ano}: primeiro vs.\ \'ultimo trimestre observado no mesmo ano $t$.
\item $^{a}$Abandono refere-se ao percentual de alunos com \texttt{V3002=1} no primeiro trimestre de $t$ que reportam n\~ao frequentar escola na \'ultima obs.\ do mesmo ano $t$. As taxas inter e intra n\~ao somam 100\% porque medem fen\^omenos em janelas temporais distintas.
\item $^{\dagger}$Anos afetados pela pandemia COVID-19 (suspens\~ao de aulas presenciais, mudan\c{c}as metodol\'ogicas na coleta da PNADC e ado\c{c}\~ao de aprova\c{c}\~ao autom\'atica em diversos estados). Interpretar com cautela.
\item \emph{N} reporta o n\'umero de observa\c{c}\~oes individuais (n\~ao ponderado) usadas no c\'alculo. As taxas s\~ao ponderadas pelo peso amostral $w_i^{P1}$ (peso PNADC longitudinalmente v\'alido). O total prom.\ + repet.\ + evas.\ + EJA $\sim$100\% no EF, mas $\sim$92\% no EM, devido a transi\c{c}\~oes do 3\textsuperscript{o} EM para o ensino superior, modalidades n\~ao classificadas e abandono escolar que se manifesta em forma de migra\c{c}\~ao para o mercado de trabalho.
\end{tablenotes}
\end{threeparttable}
\end{table}
